\chapter{Proposed Methodology}
\label{chapter_methodology}

\section{Research Approach}

The thesis will use design science with empirical evaluation. Design science is
appropriate because the research produces an artifact intended to address a
practical organizational problem and evaluates that artifact against explicit
requirements \cite{hevner2004design}. The artifact is the \system{} framework,
including its canonical data design, time-aware analytical pipeline,
constraint-aware recommendation procedure and owner-facing evidence. Empirical
study is required because a plausible architecture alone cannot establish
predictive validity, usability or business value.

The study is organized into four connected phases:

\begin{enumerate}[label=P\arabic*.]
  \item \textbf{Problem and evidence definition:} establish the target SME
  setting, research questions, literature-derived requirements, permitted
  claims and success measures.
  \item \textbf{Artifact design:} specify the operational schema, historical
  snapshots, candidate analytical methods, strategy constraints, interfaces and
  audit records.
  \item \textbf{Technical evaluation:} test data integrity, prediction,
  calibration, robustness and scalability using strictly separated development,
  public-real and primary-real evidence.
  \item \textbf{Decision and user evaluation:} assess simulated or pilot
  inventory consequences and test whether owners can understand and correctly
  use the recommendations.
\end{enumerate}

The unit of forecasting is a branch--SKU--day or branch--SKU--week, selected
according to data density. The unit of customer-risk modelling is a customer's
feature snapshot at a historical cutoff. The unit of strategy evaluation is an
eligible action evaluated under the same information and constraints as its
comparison policy.

\section{Research Questions, Hypotheses and Measures}

The research questions introduced in Chapter~\ref{chapter_introduction} are
operationalized as follows:

\begin{description}[style=nextline,leftmargin=0pt]
  \item[RQ1 / H1.] A demand model containing local calendar, price, promotion
  and available market-context features will reduce out-of-time WAPE or MASE
  relative to the otherwise identical model without those feature groups and
  relative to a seasonal-naive baseline.
  \item[RQ2 / H2.] A model trained from historical customer snapshots will
  identify future inactivity better than the event base rate, measured by
  PR-AUC, and will achieve Lift@10\% above 1.0 while retaining acceptable
  probability calibration.
  \item[RQ3 / H3.] Under identical demand sequences, costs and constraints, the
  proposed forecast-driven replenishment policy will reduce total stock-out and
  holding cost relative to the participating owner's documented current policy.
  \item[RQ4 / H4.] In counterbalanced decision tasks, an explanation condition
  will improve task correctness or confidence without a practically harmful
  increase in completion time compared with an output-only condition.
\end{description}

\begin{table}[H]
\centering
\small
\caption{Alignment of research questions, evidence and primary measures.}
\label{tab:rq_alignment}
\begin{tabularx}{\textwidth}{p{0.09\textwidth}p{0.27\textwidth}Y Y}
\toprule
\textbf{RQ} & \textbf{Unit and comparison} & \textbf{Primary measures} & \textbf{Required evidence} \\
\midrule
RQ1 & Unseen branch--SKU periods; context model versus ablations/baselines & WAPE, MASE, interval coverage & Time-stamped sales, stock, price, promotion and calendar data \\
RQ2 & Future customer windows; model versus base-rate/rule baseline & PR-AUC, Lift@10\%, Brier score, calibration & Pseudonymous customer transactions with historical cutoffs \\
RQ3 & Same demand path; proposed versus current replenishment policy & Stock-out units, holding units, total cost & Stock, lead time, order constraints and validated cost assumptions \\
RQ4 & Same decision scenarios; explanation versus output only & Correctness, time, confidence, perceived usability & Consenting target users and counterbalanced tasks \\
\bottomrule
\end{tabularx}
\end{table}

Hypotheses and primary measures are fixed before the final test is inspected.
Secondary analyses will be labelled as exploratory. A statistically detectable
difference will not be described as practically important unless its magnitude
or cost implication is material.

\section{Target Population and Sampling Plan}

The target population comprises inventory-based retail SMEs operating in
Bangladesh. To reduce uncontrolled variation, initial recruitment will focus on
three to five businesses from one broadly comparable retail vertical. The
planned analytical target is 12--24 months of records across the partners,
subject to ethics approval and data availability. A first partner with a shorter
history may be used only to test collection, mapping and data-quality rules.

Businesses will be recruited purposively because access to longitudinal invoice
and stock data requires owner agreement. Inclusion requires identifiable dates,
products and quantities and a person able to explain the current replenishment
process. Customer-level analysis additionally requires stable pseudonymous
customer identifiers and sufficient follow-up after each cutoff. A business or
task will be excluded from a particular analysis when these minimum conditions
are absent; it will not be repaired by inventing identifiers or outcomes.

\section{Data Sources and Evidence Tiers}

Evidence will be kept in separate tiers so that convenience data do not silently
change the claim being tested.

\begin{table}[H]
\centering
\small
\caption{Planned evidence tiers and claim boundaries.}
\label{tab:evidence_tiers}
\begin{tabularx}{\textwidth}{p{0.10\textwidth}p{0.29\textwidth}Y Y}
\toprule
\textbf{Tier} & \textbf{Source} & \textbf{Permitted use} & \textbf{Prohibited inference} \\
\midrule
1 & Consenting Bangladesh SME invoice, stock and operating records & Primary local forecasting, customer, strategy and usability evaluation & Generalization beyond sampled firms without further evidence \\
2a & Bangladesh WFP/DAM market-price data \cite{wfp2026prices} & Local external context and separate price-series study & SKU sales, customer behavior or firm revenue \\
2b & UCI Online Retail II, United Kingdom \cite{chen2012uci} & Real-invoice pipeline and method portability checks & Bangladesh purchasing behavior \\
2c & Public Bangladesh retailer demand series \cite{jahin2024dataset} & Local quantity-forecast procedure check & Multi-SKU, customer, price or profit claims if fields are absent \\
3 & Clearly labelled generated records & Software tests, rare cases, scale tests and privacy-safe examples & Field effectiveness or real customer behavior \\
\bottomrule
\end{tabularx}
\end{table}

Every received dataset will have a manifest containing source, owner or
repository, collection period, receipt date, checksum, license or agreement,
row count, schema version, mapping, exclusions and known limitations. Tier-3
data will never be merged into a Tier-1 test set. When augmentation is examined,
the real-only and augmented training conditions will be reported separately and
evaluated on the same untouched real test period.

\section{Data Acquisition and Variable Plan}

Primary collection will prefer exports from the business's existing point of
sale or spreadsheet workflow. A controlled template will be offered when no
consistent export exists. The minimum sales fields are timestamp, business,
branch, invoice, SKU, quantity, unit price, discount and cancellation/return
status. Recommended fields include unit cost, promotion, payment mode,
pseudonymous customer, stock-on-hand, stock movement, purchase receipt,
supplier lead time and incoming quantity.

Variables are organized by the decision time at which they become available:

\begin{table}[H]
\centering
\small
\caption{Feature availability and leakage rule.}
\label{tab:feature_availability}
\begin{tabularx}{\textwidth}{p{0.22\textwidth}p{0.30\textwidth}Y}
\toprule
\textbf{Availability class} & \textbf{Examples} & \textbf{Rule} \\
\midrule
Known before cutoff & Historical sales, current stock, recorded price, lead time, calendar & Eligible if timestamp and provenance are valid \\
Known for forecast horizon & Confirmed promotion, planned closure, calendar event & Eligible only when the same information would exist in real use \\
Observed after outcome & Later purchase, delivered satisfaction, realized return, future stock adjustment & Label or evaluation outcome only; never a predictor \\
Uncertain or revised & Unconfirmed promotion, estimated cost, manual lead time & Retain source and uncertainty; include a sensitivity analysis \\
\bottomrule
\end{tabularx}
\end{table}

\section{Ethics, Privacy and Governance}

University ethics approval and a signed data-sharing agreement will precede
primary collection. Direct identifiers such as names, phone numbers, e-mail
addresses, national identifiers, account numbers and wallet numbers are not
required for the research tables. Each business will use a separate salted
pseudonymous customer key; any re-identification mapping will remain with the
business and outside the research repository.

The agreement will state research purpose, authorized researchers, storage,
retention deadline, publication form, withdrawal and deletion procedure. Access
will follow least privilege, and published material will use aggregates or
separately reviewed de-identified examples. A retention suggestion will not
authorize contact: customer consent and the participating business's policy are
independent feasibility constraints. Model outputs will support, not replace,
owner decisions.

\section{Canonical Operational and Analytical Model}

The proposed operational schema contains Organization, Branch, Product,
Customer, Supplier, SalesOrder, SalesOrderItem, Payment, InventoryBalance,
StockMovement, PurchaseOrder, Expense, LedgerEntry, SalesReturn, Refund,
ImportBatch and AuditLog entities. The analytical layer does not train directly
on mutable operational tables. It creates versioned invoice-line, SKU-time and
customer-cutoff views with the source keys and data cutoff retained.

\begin{figure}[H]
\centering
\resizebox{\textwidth}{!}{%
\begin{tikzpicture}[
  entity/.style={draw=dugreen,rounded corners,fill=softgreen,minimum width=2.7cm,minimum height=0.9cm,align=center},
  event/.style={draw=dugold,rounded corners,fill=softorange,minimum width=2.9cm,minimum height=0.9cm,align=center},
  arrow/.style={-{Latex[length=2mm]},thick,draw=gray!75}
]
\node[entity] (org) {Organization\\Branch};
\node[entity,right=1.2cm of org] (product) {Product\\Supplier};
\node[entity,right=1.2cm of product] (customer) {Customer\\Consent};
\node[event,below=1.3cm of product] (sale) {Sale + Line Items\\Payment};
\node[event,left=1.2cm of sale] (purchase) {Purchase Order\\Receiving};
\node[event,right=1.2cm of sale] (return) {Return + Refund\\Receivable};
\node[event,below=1.3cm of sale] (stock) {Stock Movements\\Inventory Balance};
\node[entity,left=1.2cm of stock] (finance) {Expense\\Ledger};
\node[entity,right=1.2cm of stock] (audit) {Import Batch\\Audit Log};
\draw[arrow] (org) -- (sale);
\draw[arrow] (product) -- (sale);
\draw[arrow] (customer) -- (sale);
\draw[arrow] (purchase) -- (stock);
\draw[arrow] (sale) -- (stock);
\draw[arrow] (return) -- (stock);
\draw[arrow] (sale) -- (return);
\draw[arrow] (finance) -- (audit);
\draw[arrow] (stock) -- (audit);
\end{tikzpicture}}
\caption{Proposed canonical operational data model.}
\label{fig:data_model}
\end{figure}

Validation will test uniqueness of business-scoped identifiers, chronological
consistency, nonnegative physical quantities where appropriate, payment/order
reconciliation, cancellation handling and stock-movement balance. Corrections
will be recorded as transformations rather than silently overwriting the source.

\section{Formal Problem Formulation}

For business $b$, branch $l$, product $p$, customer $c$ and time $t$, let

\begin{equation}
  D_t=\{\text{Sales, Stock, Purchases, Expenses, Customers, Payments, Context}\}_{\leq t}.
\end{equation}

\subsection{Performance Evaluation}

For each KPI $K$, the system returns $K(D_t;w,g)$, where $w$ is the reporting
window and $g$ is the aggregation level. Each value retains numerator,
denominator, filters and data-quality status so that the owner can reproduce it.

\subsection{Demand Prediction}

\begin{equation}
\hat{y}_{b,l,p,t+h}=f\!\left(X^{\mathrm{history}}_{t},X^{\mathrm{price}}_{t},
X^{\mathrm{promotion}}_{t},X^{\mathrm{calendar}}_{t:t+h},
X^{\mathrm{market}}_{t}\right),
\label{eq:demand}
\end{equation}
where $\hat{y}$ is expected demand at horizon $h$. Quantile forecasts or
calibrated residual intervals will represent uncertainty. When product-level
forecasts are aggregated, reconciliation will be considered so SKU, category,
branch and business totals remain coherent \cite{hyndman2011hierarchical}.

\subsection{Future Customer Inactivity}

\begin{equation}
q_{c,t}=P\!\left(\text{no purchase in }(t,t+H]\mid X_{c,t}\right),
\label{eq:churn}
\end{equation}
where $X_{c,t}$ contains only records available by cutoff $t$. Horizon $H$ will
be fixed from business purchase cycles before final evaluation.

\subsection{Strategy Optimization}

For an eligible action $a\in A_t$,

\begin{equation}
U(a\mid D_t)=E[\Delta\mathrm{Profit}(a)]-\mathrm{ActionCost}(a)
-\lambda\,\mathrm{Risk}(a),
\label{eq:utility}
\end{equation}
subject to budget, storage, minimum order quantity, supplier lead time,
current/incoming stock and contact consent. This is a transparent operational
objective inspired by the predictive-to-prescriptive transition
\cite{bertsimas2020prescriptive}; its terms will be estimated and sensitivity
tested rather than presented as known constants.

\section{Preprocessing and Historical Snapshots}

The processing order is fixed: validate source and schema; remove or mark exact
duplicates; resolve cancellations and returns; construct the split; fit
imputation, encoding and scaling on training data; build features; train; and
evaluate once on the untouched test. The order prevents statistics from future
periods entering training \cite{kaufman2012leakage}.

Demand features may include lags, rolling location and variability, day of week,
month, festival distance, price, discount, known promotion, stock-out indicator
and external market context. Zero demand caused by no sale is distinguished,
when possible, from zero observed sales caused by no stock. Intermittency will
be measured before choosing temporal granularity and candidate methods.

Customer snapshots will include RFM, tenure, order-value statistics, category
breadth, purchase interval, discount behavior and pre-cutoff return rate. The
label window follows the observation window, with a purge interval when feature
construction or delayed events could overlap it. No customer identifier appears
on both sides of an entity-generalization test.

\section{Candidate Methods and Selection Rules}

Method selection follows a baseline ladder, and eligibility depends on evidence
volume and task structure.

\begin{itemize}
  \item \textbf{Performance:} auditable KPI definitions and rule-based anomaly
  flags establish the descriptive baseline.
  \item \textbf{Demand:} last-value, seasonal-naive, moving-average or
  exponential-smoothing baselines precede ARIMA/SARIMA, Croston-family methods,
  gradient-boosted lag models and, only with adequate sequences, LSTM. Croston
  and the Syntetos--Boylan correction address intermittent demand
  \cite{croston1972intermittent,syntetos2005intermittent}.
  \item \textbf{Inactivity:} a recency or base-rate rule is compared with
  Logistic Regression, Random Forest \cite{breiman2001rf} and XGBoost
  \cite{chen2016xgboost}. Class weighting and decision threshold are selected on
  development data.
  \item \textbf{Segmentation:} scaled RFM features are studied with K-Means.
  Candidate $K$ values are judged by inertia, silhouette, resampling stability
  and business interpretability.
  \item \textbf{Explanation:} coefficient/rule explanations are preferred for
  simple models; SHAP supplies model-behavior attributions for tree models
  \cite{lundberg2017shap,lundberg2020trees}. Displayed reasons will remain
  selective and contextual \cite{miller2019explanation}.
\end{itemize}

A complex candidate will be rejected when it lacks sufficient training history,
cannot beat the relevant baseline on development data, produces unstable
probabilities, or adds operational cost without material value. This rule keeps
the proposed system useful when the scientifically honest answer is a simple
forecast or business rule.

\section{The \system{} Sequential Procedure}

\begin{algorithm}[H]
\caption{\system{} Strategy Recommendation}
\label{alg:bsmart}
\KwIn{Data available at cutoff $D_t$, horizon $H$, constraints $C_t$, validated model/baseline versions $M$, candidate actions $A$}
\KwOut{Ranked feasible recommendations $R_t$}
Verify source, schema, timestamps and organization ownership\;
Create the historical feature snapshot $X_t$\;
Apply transformations fitted on the appropriate training evidence\;
Estimate demand with uncertainty and customer inactivity probabilities\;
Calculate reproducible KPIs and data-quality warnings\;
Generate candidate actions with source and decision-time metadata\;
\ForEach{$a\in A$}{
  Check feasibility under $C_t$\;
  Estimate benefit, cost and uncertainty using declared assumptions\;
  $U(a)\leftarrow E[\Delta\mathrm{Profit}(a)]-\mathrm{Cost}(a)-\lambda\mathrm{Risk}(a)$\;
}
Remove infeasible actions and rank the remainder by $U(a)$\;
Attach concise explanation, uncertainty, cutoff and evidence source\;
Return top-$K$ actions; record owner response and later observable outcome\;
\end{algorithm}

For $N$ invoice lines, $d$ validated fields and $P$ candidate actions, source
validation and feature aggregation are approximately $O(Nd)$ and action ranking
is $O(P\log P)$. Training complexity depends on the selected model and is
measured empirically. Serving reuses versioned artifacts; it must not retrain in
the request path.

\section{Proposed Architecture and Information Flow}

\begin{figure}[H]
\centering
\resizebox{0.98\textwidth}{!}{%
\begin{tikzpicture}[
  block/.style={draw,rounded corners,minimum width=3.0cm,minimum height=1.05cm,align=center,font=\small},
  arrow/.style={-{Latex[length=2.2mm]},thick},node distance=0.75cm
]
\node[block,fill=softgreen] (ops) {Business Operations\\sales, stock, purchases, returns};
\node[block,fill=softblue,right=of ops] (canon) {Canonical Evidence\\validation, provenance, snapshots};
\node[block,fill=softorange,right=of canon] (ai) {AI Analytics\\forecast, risk, segments, explanation};
\node[block,fill=softgreen,right=of ai] (strategy) {Strategy Layer\\constraints, utility, ranking};
\node[block,fill=softblue,below=1.0cm of strategy] (ui) {Owner Interface\\reason, uncertainty, action};
\node[block,fill=softgray,left=of ui] (outcome) {Outcome Record\\response and observed result};
\node[block,fill=softgray,left=of outcome] (monitor) {Monitoring\\quality, drift, calibration};
\draw[arrow] (ops) -- (canon);
\draw[arrow] (canon) -- (ai);
\draw[arrow] (ai) -- (strategy);
\draw[arrow] (strategy) -- (ui);
\draw[arrow] (ui) -- (outcome);
\draw[arrow] (outcome) -- (monitor);
\draw[arrow] (monitor.west) -| ($(ops.south)+(0,-0.15)$) -- (ops.south);
\end{tikzpicture}}
\caption{Proposed \system{} architecture and accountable feedback loop.}
\label{fig:bsmart_architecture}
\end{figure}

The information flow is intentionally one-directional at prediction time:
future outcomes cannot alter the snapshot that produced an earlier result.
Outcomes enter later through the monitoring path and become eligible training
evidence only in a new, versioned cycle. This separation makes it possible to
reproduce what the system knew when an action was proposed.

\section{Evaluation Protocol}

\subsection{Forecasting Evaluation}

All forecasting splits preserve time. Candidate methods are selected through
rolling-origin or expanding-window validation, followed by one evaluation on a
final untouched period. For actual values $y_i$ and forecasts $\hat y_i$,

\begin{align}
\mathrm{MAE} &= \frac{1}{n}\sum_{i=1}^{n}|y_i-\hat y_i|,\\
\mathrm{RMSE} &= \sqrt{\frac{1}{n}\sum_{i=1}^{n}(y_i-\hat y_i)^2},\\
\mathrm{WAPE} &= \frac{\sum_i|y_i-\hat y_i|}{\sum_i|y_i|}.
\end{align}

MASE will support comparison across scales and avoid some defects of percentage
errors \cite{hyndman2006accuracy}. Prediction-interval coverage and width or a
proper scoring rule will evaluate uncertainty \cite{gneiting2007scoring}.
Results will be stratified by demand volume and intermittency rather than hiding
failure on low-volume items in one aggregate.

\subsection{Customer and Segmentation Evaluation}

Inactivity models will report outcome prevalence, PR-AUC, ROC-AUC, Lift@10\%,
Brier score and reliability plots. The emphasis on PR measures follows their
value under imbalance \cite{saito2015pr}; calibration is assessed separately
because a useful ranking can still produce misleading probabilities
\cite{niculescu2005calibration}. Thresholds and contact fractions are chosen on
validation evidence using declared costs, not adjusted after viewing the test.

Clustering evaluation will report inertia, silhouette, cluster stability,
minimum cluster size and RFM profiles. The final interpretation will be reviewed
with domain participants. A low-separation result will be reported as such; the
existence of clusters will not be treated as proof of natural customer types.

\subsection{Ablation, Robustness and Statistical Analysis}

Planned ablations remove one feature family at a time: local calendar,
price/promotion, market context and customer-history extensions. Additional
comparisons examine global versus business-specific models and real-only versus
clearly controlled augmentation. All paired comparisons use the same test
instances or forecast origins.

Uncertainty will be summarized with paired block bootstrap intervals across
time or business units, as appropriate. Model differences will be reported with
effect size and interval, not only a significance label. Sensitivity analysis
will vary lead time, shortage cost, holding cost, action cost and risk penalty in
Equation~\ref{eq:utility}.

\subsection{Decision, Usability and Scalability Evaluation}

Replenishment policies will face identical observed or replayed demand paths.
The comparison records stock-out units, average inventory, holding cost,
shortage cost and total cost. Retention ranking will use expected value only
when response benefit, contact cost and consent are available; otherwise it will
remain a prioritization study, not a profit claim \cite{verbeke2012profit}.

Owner evaluation will use realistic tasks in a counterbalanced within-subject
design when the sample permits. Measures include task correctness, completion
time, error count, confidence and perceived usability. Participants may explain
their decision in their own words, allowing the study to detect plausible but
misunderstood explanations.

Scalability tests will use generated data only and will remain separate from
accuracy results. Planned sizes are 100, 500, 1,000, 10,000, 100,000 and
1,000,000 transaction lines. Measures include import and preprocessing time,
training time, prediction latency, API throughput, peak memory, artifact size
and representative query latency. Any association-rule support experiment will
be reported only after the corresponding algorithm exists and has been run.

\subsection{Monitoring Plan}

Monitoring will distinguish schema or data-quality failure, input shift,
calibration change and outcome-performance change. Alerts will not automatically
authorize retraining. A new model must repeat the validation protocol and retain
its training window, features, parameters and comparison result. This follows
the concept-drift literature's emphasis on both adaptation and appropriate
evaluation \cite{gama2014drift}.

\section{Threats to Validity and Planned Controls}

\begin{table}[H]
\centering
\small
\caption{Principal validity threats and planned controls.}
\label{tab:validity_controls}
\begin{tabularx}{\textwidth}{p{0.22\textwidth}p{0.30\textwidth}Y}
\toprule
\textbf{Threat} & \textbf{How it may distort evidence} & \textbf{Planned control} \\
\midrule
Selection bias & Cooperative firms may be more digitally mature & Describe recruitment and firms; restrict generalization \\
Leakage & Future or target-derived data inflate performance & Availability table, cutoff snapshots, grouped/temporal split \\
Stock-out censoring & Observed sales understate unconstrained demand & Retain stock state; stratify or flag censored periods \\
Small/short series & Complex models overfit & Eligibility rule, simple baselines, rolling validation \\
Dataset mismatch & Foreign/public tasks are mistaken for local evidence & Evidence tiers and separate result tables \\
Uncertain costs & Utility ranking depends on weak assumptions & Owner-validated ranges and sensitivity analysis \\
Novelty/learning effect & Interface condition benefits from presentation order & Counterbalance tasks and collect task-level timing \\
Concept drift & Festival, price or behavior relations change & Time-based tests and post-deployment monitoring plan \\
\bottomrule
\end{tabularx}
\end{table}

\section{Two-Semester Work Plan}

The project period is February 2026--February 2027. The first semester defines
the problem, evidence base and proposed artifact; the second emphasizes primary
data, implementation completion, controlled evaluation and consolidation. The
schedule in Table~\ref{tab:gantt} is dependency-aware: the final model study
cannot precede ethics/data readiness, and the final conclusions cannot precede
the frozen evaluation.

\begin{landscape}
\begin{table}[H]
\centering
\scriptsize
\setlength{\tabcolsep}{3.2pt}
\renewcommand{\arraystretch}{1.35}
\caption{Project timeline from February 2026 to February 2027.}
\label{tab:gantt}
\begin{tabular}{p{0.24\linewidth}*{13}{c}}
\toprule
\textbf{Work package} & \rotatebox{90}{Feb 26} & \rotatebox{90}{Mar 26} & \rotatebox{90}{Apr 26} & \rotatebox{90}{May 26} & \rotatebox{90}{Jun 26} & \rotatebox{90}{Jul 26} & \rotatebox{90}{Aug 26} & \rotatebox{90}{Sep 26} & \rotatebox{90}{Oct 26} & \rotatebox{90}{Nov 26} & \rotatebox{90}{Dec 26} & \rotatebox{90}{Jan 27} & \rotatebox{90}{Feb 27} \\
\midrule
WP1: Background study and problem formulation & \cellcolor{softgreen} & \cellcolor{softgreen} & \cellcolor{softgreen} & & & & & & & & & & \\
WP2: Requirements, ethics and data protocol & & & \cellcolor{softgreen} & \cellcolor{softgreen} & \cellcolor{softgreen} & \cellcolor{softgreen} & & & & & & & \\
WP3: Framework, schema and evaluation design & & & & \cellcolor{softgreen} & \cellcolor{softgreen} & \cellcolor{softgreen} & \cellcolor{softgreen} & \cellcolor{softgreen} & & & & & \\
WP4: Core artifact implementation & & & & & & \cellcolor{softblue} & \cellcolor{softblue} & \cellcolor{softblue} & \cellcolor{softblue} & & & & \\
WP5: Partner onboarding and primary data & & & & & & & & \cellcolor{softblue} & \cellcolor{softblue} & \cellcolor{softblue} & & & \\
WP6: Model development and internal validation & & & & & & & & & \cellcolor{softblue} & \cellcolor{softblue} & \cellcolor{softblue} & & \\
WP7: Frozen tests, ablations and scale study & & & & & & & & & & & \cellcolor{softorange} & \cellcolor{softorange} & \\
WP8: Usability and decision-impact study & & & & & & & & & & & & \cellcolor{softorange} & \cellcolor{softorange} \\
WP9: Final writing, audit and defence preparation & & & & & & & & & & & \cellcolor{softorange} & \cellcolor{softorange} & \cellcolor{softorange} \\
\bottomrule
\end{tabular}

\vspace{0.3cm}
\textit{Legend:} \colorbox{softgreen}{Phase-I research foundation}\quad
\colorbox{softblue}{build and data preparation}\quad
\colorbox{softorange}{evaluation and consolidation}.
\end{table}
\end{landscape}

\begin{table}[H]
\centering
\small
\caption{Milestones, exit criteria and contingency.}
\label{tab:milestones}
\begin{tabularx}{\textwidth}{p{0.17\textwidth}p{0.18\textwidth}Y Y}
\toprule
\textbf{Milestone} & \textbf{Target} & \textbf{Exit criterion} & \textbf{Contingency} \\
\midrule
M1: Initial report & September 2026 & Problem, verified background, proposed method and approved plan & Revise scope/RQs before model lock \\
M2: Data readiness & November 2026 & Ethics/agreement complete; manifest and quality report for eligible primary data & Narrow the primary task; retain public data only as external procedure checks \\
M3: Model freeze & December 2026 & Baselines, features, splits, metrics and seeds fixed before final test & Prefer simpler eligible models when volume is inadequate \\
M4: Evaluation freeze & January 2027 & Untouched tests, ablations, intervals, decision and scale analyses complete & Report inconclusive/negative findings; do not retune on test \\
M5: Final thesis & February 2027 & Chapters complete; citations, claims, figures and reproducibility artifacts audited & Remove unsupported claims and document remaining limitations \\
\bottomrule
\end{tabularx}
\end{table}

The critical path runs through ethics and primary-data readiness (M2), model
freeze (M3) and evaluation freeze (M4). Weekly progress review will track a
deliverable, its evidence and its next dependency. Schedule pressure will not be
handled by converting a public or generated-data result into a primary local
claim.

\section{Expected Methodological Outputs}

At completion, the methodology is expected to yield a documented canonical
schema and data dictionary, partner data protocol, versioned feature snapshots,
baseline and candidate model specifications, a constraint-aware recommendation
procedure, a reproducible evaluation package and an owner-study instrument.
Implementation detail, observed numerical results and conclusions will be added
to the reserved chapters only after the corresponding milestones are satisfied.
